\section{Conclusion}
In this paper, we present a reward design tailored for GRPO training on tool use tasks. Empirically, our model trained from scratch using GRPO consistently outperforms both SFT-based and SFT-initialized RL baselines, as well as models trained with alternative RL algorithms, across a variety of held-out tool use benchmarks. Furthermore, we demonstrate that our model generalizes well to QA settings, exhibiting robust multi-turn interactions, emergent proactiveness, and metacognitive behaviors, all of which are key traits for efficient and adaptable tool use, lying at the core of foundational agent capabilities.
Our in-depth analysis of reward types, scaling strategies, granularity, and temporal dynamics provides further insights into how reward shaping influences learning and behavior. We hope these findings serve as a roadmap for future work in applying reinforcement learning to tool use. Ultimately, we envision that \textit{reward is all tool learning needs}, and that RL offers a powerful path toward generalizable and creative agent behavior.


% \section*{Limitations}
% Our study focuses on three key domains where LLMs explicitly struggle—Math, Intention, and Time—building on insights from existing literature. However, LLMs also face challenges in areas such as long-tail knowledge and domain-specific expertise, where external resources are essential. Expanding SMART-ER to these domains could further refine model self-awareness and improve calibration in knowledge boundary, complementing the strong OOD  performance that SMARTAgent has already demonstrated. Additionally, while we evaluate our approach on two major model families, extending our analysis to a broader range of architectures, including Qwen, DeepSeek, and varying model sizes, could further validate and enhance the generalizability of our findings.
